% Phần khai báo dùng chung cho báo cáo (biên dịch bằng XeLaTeX/Tectonic).
% Cỡ chữ thân bài 13pt (lớp report chỉ có sẵn 10, 11, 12pt nên phải đặt qua gói này).
\usepackage[fontsize=13pt]{scrextend}
% Lề theo template của Khoa: trái 3,5 cm, phải 2 cm, trên 3,5 cm, dưới 3 cm; số trang ở đầu trang, giữa.
\usepackage[left=3.5cm,right=2cm,top=3.5cm,bottom=3cm,headheight=16pt,headsep=1.5cm,footskip=11mm]{geometry}
\usepackage{fontspec}
\setmainfont{Times New Roman}
\setsansfont{Arial}
\usepackage{amsmath,amssymb,graphicx,booktabs,longtable,tabularx,array,multirow}
\usepackage{enumitem,fancyhdr,titlesec,titletoc,caption,needspace,xurl,float,setspace,indentfirst}
\usepackage[hidelinks,unicode,pdftitle={A Book Advisory Assistant Based on Multiple Agents},pdfauthor={Hoàng Xuân Thành, Trần Gia Lạc}]{hyperref}
\usepackage{tikz}
\usepackage{adjustbox}
\usetikzlibrary{positioning,arrows.meta,patterns,calc,shapes.geometric}
\usepackage{pgfplots}
\pgfplotsset{compat=1.17}
\hypersetup{bookmarksnumbered=true}

% ---------- English names ----------
\renewcommand{\figurename}{Figure}
\renewcommand{\tablename}{Table}
\renewcommand{\contentsname}{Contents}
\renewcommand{\listfigurename}{List of figures}
\renewcommand{\listtablename}{List of tables}
\renewcommand{\chaptername}{Chapter}
\renewcommand{\appendixname}{Appendix}

% ---------- Đoạn văn, danh sách ----------
% Template: giãn dòng 1,5 lines của Word (Times New Roman 13pt ~ 22,4pt), thụt đầu dòng 1,27 cm,
% không cách đoạn.
\setstretch{1.44}
\setlength{\parindent}{1.27cm}
\setlength{\parskip}{0pt}
\setlist{leftmargin=1.2cm,itemsep=2pt,topsep=3pt,parsep=1pt}
\setlength{\emergencystretch}{3em}
\widowpenalty=10000
\clubpenalty=10000
% Không ngắt trang giữa câu dẫn và danh sách đi ngay sau nó
\makeatletter\@beginparpenalty=10000\makeatother
\raggedbottom

% ---------- Bảng ----------
\renewcommand{\arraystretch}{1.25}
\newcolumntype{Y}{>{\raggedright\arraybackslash}X}
\newcolumntype{P}[1]{>{\raggedright\arraybackslash}p{#1}}
\newcolumntype{R}[1]{>{\raggedleft\arraybackslash}p{#1}}
\newcolumntype{C}[1]{>{\centering\arraybackslash}p{#1}}
% Template: "Hình 2.1: ..." và "Bảng 4.1: ...", chữ thường, căn giữa.
\captionsetup{font=normalsize,labelsep=colon,justification=centering}
\captionsetup[table]{position=top,skip=6pt}
\captionsetup[figure]{position=bottom,skip=8pt}
% Bảng và hình được đặt gần chỗ nhắc tới; nếu không đủ chỗ thì chuyển lên đầu trang sau
% và phần chữ tiếp tục lấp đầy trang hiện tại (tránh để trống nửa trang).
\renewcommand{\topfraction}{0.9}
\renewcommand{\bottomfraction}{0.7}
\renewcommand{\textfraction}{0.08}
\renewcommand{\floatpagefraction}{0.8}
\setcounter{topnumber}{3}
\setcounter{bottomnumber}{2}
\setcounter{totalnumber}{4}
\setlength{\textfloatsep}{14pt plus 2pt minus 2pt}
\setlength{\intextsep}{12pt plus 2pt minus 2pt}
% Trang chỉ có bảng/hình: xếp từ đầu trang xuống, không dàn đều theo chiều dọc
\makeatletter
\setlength{\@fptop}{0pt}
\setlength{\@fpsep}{22pt}
\makeatother

% ---------- Tiêu đề chương, mục ----------
% Template: "CHƯƠNG 1. TÊN CHƯƠNG" 16pt đậm, căn trái; mục cấp 1 (1.1) 14pt đậm;
% cấp 2 (1.1.1) 14pt đậm nghiêng; cấp 3 (1.1.1.1) 14pt thường.
% Tiêu đề không đánh số (Lời cảm ơn, Mục lục, ...) 16pt đậm, căn giữa.
\newcommand{\tienochuong}{CHAPTER}
\titleformat{\chapter}[hang]{\normalfont\bfseries\fontsize{16}{21}\selectfont\raggedright\hyphenpenalty=10000}{\tienochuong~\thechapter.}{0.5em}{\MakeUppercase}
\titleformat{name=\chapter,numberless}[block]{\normalfont\bfseries\centering\fontsize{16}{21}\selectfont}{}{0pt}{\MakeUppercase}
\titlespacing*{\chapter}{0pt}{-20pt}{14pt}
\titleformat{\section}{\normalfont\bfseries\fontsize{14}{18}\selectfont}{\thesection}{0.6em}{}
\titleformat{\subsection}{\normalfont\bfseries\itshape\fontsize{14}{18}\selectfont}{\thesubsection}{0.6em}{}
\titleformat{\subsubsection}{\normalfont\fontsize{14}{18}\selectfont}{\thesubsubsection}{0.6em}{}
\titlespacing*{\section}{0pt}{12pt}{4pt}
\titlespacing*{\subsection}{0pt}{9pt}{3pt}
\titlespacing*{\subsubsection}{0pt}{7pt}{2pt}
% Không để tiêu đề mục rơi xuống cuối trang mà gần như không có nội dung đi kèm
\usepackage{etoolbox}
\pretocmd{\section}{\Needspace{6\baselineskip}}{}{}
\pretocmd{\subsection}{\Needspace{4\baselineskip}}{}{}
\setcounter{secnumdepth}{3}
\setcounter{tocdepth}{2}

% ---------- Mục lục, danh mục hình và bảng (theo template) ----------
% Dòng chương: "CHƯƠNG 1. TÊN CHƯƠNG" đậm; mục 1.1 chữ thường; mục 1.1.1 nghiêng.
\newcommand{\tienomucluc}{CHAPTER}
\titlecontents{chapter}[0pt]{\addvspace{2pt}\bfseries}
  {\tienomucluc~\thecontentslabel.\ \MakeUppercase}{\MakeUppercase}
  {\titlerule*[0.6pc]{.}\contentspage}
\titlecontents{section}[0.4cm]{}{\thecontentslabel\ \ }{}{\titlerule*[0.6pc]{.}\contentspage}
\titlecontents{subsection}[0.8cm]{\itshape}{\thecontentslabel\ \ }{}{\titlerule*[0.6pc]{.}\contentspage}
\titlecontents{figure}[0pt]{}{\figurename~\thecontentslabel:\ }{}{\titlerule*[0.6pc]{.}\contentspage}
\titlecontents{table}[0pt]{}{\tablename~\thecontentslabel:\ }{}{\titlerule*[0.6pc]{.}\contentspage}

% ---------- Đầu trang, chân trang ----------
% Template: chỉ có số trang ở giữa đầu trang, không có dòng tiêu đề.
\pagestyle{fancy}
\fancyhf{}
\fancyhead[C]{\thepage}
\renewcommand{\headrulewidth}{0pt}
\fancypagestyle{plain}{\fancyhf{}\fancyhead[C]{\thepage}\renewcommand{\headrulewidth}{0pt}}

% ---------- Lệnh tiện dụng ----------
\newcommand{\src}[1]{\hyperlink{src-#1}{[#1]}}
% \hinh{tệp trong figures/}{chú thích}{nhãn}
% Hình rộng hơn vùng chữ (lề trái 3,5 cm theo template) được thu nhỏ vừa khít.
\newcommand{\hinh}[3]{\begin{figure}[!htbp]\centering{\cochuhinh\begin{adjustbox}{max width=\linewidth}\input{figures/#1.tex}\end{adjustbox}}\caption{#2}\label{#3}\end{figure}}
% Cỡ chữ cố định bên trong hình: không phụ thuộc cỡ chữ thân bài, để kích thước các ô đã căn không bị lệch.
\newcommand{\cochuhinh}{%
  \renewcommand{\normalsize}{\fontsize{11}{13.2}\selectfont}%
  \renewcommand{\small}{\fontsize{10}{12.2}\selectfont}%
  \renewcommand{\footnotesize}{\fontsize{9}{10.7}\selectfont}%
  \renewcommand{\scriptsize}{\fontsize{8}{9.1}\selectfont}%
  \normalsize}
% Cỡ chữ trong bảng
\newcommand{\cobang}{\fontsize{11}{13.2}\selectfont}
\urlstyle{same}
\newcommand{\pp}{pp}

% ---------- Kiểu vẽ hình: chỉ đen, trắng, xám; nét đậm để in rõ ----------
\tikzset{
  baocao/.style={font=\sffamily\footnotesize, line width=0.9pt, >={Stealth[length=2.8mm,width=2.2mm]}},
  hop/.style={draw=black, fill=white, rounded corners=2pt, align=center, inner sep=5pt, minimum height=9mm},
  hopxam/.style={hop, fill=black!14},
  hopden/.style={hop, fill=black!85, text=white},
  khung/.style={draw=black, dashed, line width=0.8pt, rounded corners=3pt},
  tenkhung/.style={font=\sffamily\footnotesize\bfseries, anchor=north west, inner sep=3pt},
  ghichu/.style={font=\sffamily\footnotesize, align=center, inner sep=2pt},
  mui/.style={->, line width=1pt},
  haichieu/.style={<->, line width=1pt},
  so/.style={circle, draw=black, fill=black, text=white, font=\sffamily\scriptsize\bfseries, inner sep=1.2pt, minimum size=4.6mm},
}

% ---------- Kiểu biểu đồ cột: phân biệt bằng độ đậm và nét gạch, không cần màu ----------
\pgfplotsset{
  bieudo/.style={
    ybar, ymin=0, ymax=118, ytick={0,20,40,60,80,100},
    axis lines*=left, axis line style={line width=0.8pt},
    ymajorgrids, grid style={black!30, line width=0.3pt},
    tick style={black, line width=0.6pt},
    tick label style={font=\sffamily\footnotesize, /pgf/number format/assume math mode=true},
    label style={font=\sffamily\footnotesize},
    xtick=data, xtick style={draw=none},
    x tick label style={font=\sffamily\footnotesize, align=center, text width=2.2cm},
    enlarge x limits=0.11,
    nodes near coords,
    every node near coord/.append style={font=\sffamily\scriptsize, inner sep=1.5pt,
      /pgf/number format/use period, /pgf/number format/fixed, /pgf/number format/fixed zerofill, /pgf/number format/precision=1, /pgf/number format/assume math mode=true},
    legend style={font=\sffamily\footnotesize, draw=none, at={(0.5,-0.3)}, anchor=north, legend columns=4,
      /tikz/every even column/.append style={column sep=4mm}},
    legend cell align=left,
    area legend,
  },
  % nhãn số viết dọc phía trên cột (dùng khi cột hẹp)
  nhandoc/.style={every node near coord/.append style={rotate=90, anchor=west}},
  cotA/.style={fill=white, draw=black, line width=0.7pt},
  cotB/.style={pattern=north east lines, pattern color=black, draw=black, line width=0.7pt},
  cotC/.style={fill=black!45, draw=black, line width=0.7pt},
  cotD/.style={fill=black, draw=black, line width=0.7pt},
}
